% GENERATED FILE. Edit canonical lessons, then run npm run generate:books.
% canonical-source-hash: fnv1a64:fa96819041ccd5a6
% canonical-lessons: RU-C188-poleznyy, RU-C188-vrednyy, RU-C188-nuzhnyy, RU-C188-vozmozhnyy, RU-C188-slozhnyy

\chapter{Useful, Harmful, Necessary, Possible, Complicated}
\label{ch:ru-useful-harmful-necessary-possible-complicated}
\hlchaptermodality{188}

\begin{chapteropening}
\textbf{By the end of this chapter:} \emph{I can say useful, harmful, necessary, possible and complicated.}

Complete the last lesson of chapter 188: I can say useful, harmful, necessary, possible and complicated.
\end{chapteropening}

\section[\texorpdfstring{\ru{полезный} (poléznyy)}{полезный (poléznyy)}]{\ru{полезный} (poléznyy) — useful}

\label{lesson:RU-C188-poleznyy}



\begin{quote}
Before the new one: say the Russian for square, then the Russian for sharp; spicy.
\end{quote}

\subsection*{You'll want to know: \ru{полезный}}
\textbf{\ru{полезный}} (\emph{poléznyy}) — \textquotedblleft{}useful\textquotedblright{}.

\begin{grammarlens}[title={What you've built}]
One more word for this chapter.
\end{grammarlens}

\subsection*{Your turn}
\begin{itemize}
\raggedright
  \item \emph{Say it:} \emph{poléznyy}
  \item \emph{Say it:} \emph{poléznyy}, once more
  \item \emph{Say it:} say \emph{poléznyy}
  \item \emph{Recall:} say the Russian for square, then the Russian for sharp; spicy, then say \emph{poléznyy} again
\end{itemize}

\begin{tcolorbox}[breakable,colback=teal!4,colframe=teal!35!black,title={Before you move on}]
What does \ru{полезный} mean? (\textquotedblleft{}Useful\textquotedblright{}.) Say it once more.
\end{tcolorbox}
\section[\texorpdfstring{\ru{вредный} (vrédnyy)}{вредный (vrédnyy)}]{\ru{вредный} (vrédnyy) — harmful}

\label{lesson:RU-C188-vrednyy}



\begin{quote}
Before the new one: say the Russian for sharp; spicy, then the Russian for useful.
\end{quote}

\subsection*{You'll want to know: \ru{вредный}}
\textbf{\ru{вредный}} (\emph{vrédnyy}) — \textquotedblleft{}harmful\textquotedblright{}.

\begin{grammarlens}[title={What you've built}]
One more word for this chapter.
\end{grammarlens}

\subsection*{Your turn}
\begin{itemize}
\raggedright
  \item \emph{Say it:} \emph{vrédnyy}
  \item \emph{Say it:} \emph{vrédnyy}, once more
  \item \emph{Say it:} say \emph{vrédnyy}
  \item \emph{Recall:} say the Russian for sharp; spicy, then the Russian for useful, then say \emph{vrédnyy} again
\end{itemize}

\begin{tcolorbox}[breakable,colback=teal!4,colframe=teal!35!black,title={Before you move on}]
What does \ru{вредный} mean? (\textquotedblleft{}Harmful\textquotedblright{}.) Say it once more.
\end{tcolorbox}
\section[\texorpdfstring{\ru{нужный} (núzhnyy)}{нужный (núzhnyy)}]{\ru{нужный} (núzhnyy) — necessary}

\label{lesson:RU-C188-nuzhnyy}



\begin{quote}
Before the new one: say the Russian for useful, then the Russian for harmful.
\end{quote}

\subsection*{You'll want to know: \ru{нужный}}
\textbf{\ru{нужный}} (\emph{núzhnyy}) — \textquotedblleft{}necessary\textquotedblright{}.

\begin{grammarlens}[title={What you've built}]
One more word for this chapter.
\end{grammarlens}

\subsection*{Your turn}
\begin{itemize}
\raggedright
  \item \emph{Say it:} \emph{núzhnyy}
  \item \emph{Say it:} \emph{núzhnyy}, once more
  \item \emph{Say it:} say \emph{núzhnyy}
  \item \emph{Recall:} say the Russian for useful, then the Russian for harmful, then say \emph{núzhnyy} again
\end{itemize}

\begin{tcolorbox}[breakable,colback=teal!4,colframe=teal!35!black,title={Before you move on}]
What does \ru{нужный} mean? (\textquotedblleft{}Necessary\textquotedblright{}.) Say it once more.
\end{tcolorbox}
\section[\texorpdfstring{\ru{возможный} (vozmózhnyy)}{возможный (vozmózhnyy)}]{\ru{возможный} (vozmózhnyy) — possible}

\label{lesson:RU-C188-vozmozhnyy}



\begin{quote}
Before the new one: say the Russian for harmful, then the Russian for necessary.
\end{quote}

\subsection*{You'll want to know: \ru{возможный}}
\textbf{\ru{возможный}} (\emph{vozmózhnyy}) — \textquotedblleft{}possible\textquotedblright{}.

\begin{grammarlens}[title={What you've built}]
One more word for this chapter.
\end{grammarlens}

\subsection*{Your turn}
\begin{itemize}
\raggedright
  \item \emph{Say it:} \emph{vozmózhnyy}
  \item \emph{Say it:} \emph{vozmózhnyy}, once more
  \item \emph{Say it:} say \emph{vozmózhnyy}
  \item \emph{Recall:} say the Russian for harmful, then the Russian for necessary, then say \emph{vozmózhnyy} again
\end{itemize}

\begin{tcolorbox}[breakable,colback=teal!4,colframe=teal!35!black,title={Before you move on}]
What does \ru{возможный} mean? (\textquotedblleft{}Possible\textquotedblright{}.) Say it once more.
\end{tcolorbox}
\section[\texorpdfstring{\ru{сложный} (slózhnyy)}{сложный (slózhnyy)}]{\ru{сложный} (slózhnyy) — complicated}

\label{lesson:RU-C188-slozhnyy}



\begin{quote}
Before the new one: say the Russian for necessary, then the Russian for possible.
\end{quote}

\subsection*{You'll want to know: \ru{сложный}}
\textbf{\ru{сложный}} (\emph{slózhnyy}) — \textquotedblleft{}complicated\textquotedblright{}.

\begin{grammarlens}[title={What you've built}]
One more word for this chapter.
\end{grammarlens}

\subsection*{Your turn}
\begin{itemize}
\raggedright
  \item \emph{Say it:} \emph{slózhnyy}
  \item \emph{Say it:} \emph{slózhnyy}, once more
  \item \emph{Say it:} say \emph{slózhnyy}
  \item \emph{Recall:} say the Russian for necessary, then the Russian for possible, then say \emph{slózhnyy} again
\end{itemize}

\begin{tcolorbox}[breakable,colback=teal!4,colframe=teal!35!black,title={Before you move on}]
What does \ru{сложный} mean? (\textquotedblleft{}Complicated\textquotedblright{}.) Say it once more.
\end{tcolorbox}
